
\chapter*{Platonic Introduction to Volume II: the Swiss-cheese and bulk--boundary--line refoundation}
\addcontentsline{toc}{chapter}{Platonic Introduction to Volume II}

\section*{What changes in Volume II}

If Volume~I built the logarithm and its modular tower, Volume~II asks what that logarithm actually does in three dimensions.  The answer is not ``it produces a few extra operations.''  The answer is: it becomes a full holomorphic-topological organism.  The bar complex is re-read as a Swiss-cheese object; its cohomological shadow becomes a Poisson vertex algebra; its dual open sector becomes a category of lines; its interactions are governed by brace operations, spectral $R$-matrices, and bulk--boundary--line triangles; and the full homotopy package becomes the natural language rather than a background correction.

This volume therefore carries the strongest form of the upgrade requested in your prompt.  We no longer treat the convolution dg Lie algebra as the final deformation object and the modular package as an appendix.  Instead, every chapter-scale unit is introduced from the standpoint that strict dg structures are the visible spine of a larger homotopy-coherent world.  In the 3d holomorphic-topological setting this is exactly what physics forces on us: BRST anomaly brackets, defect operations, line OPE's, and boundary transfer are all higher operations before they are strict algebras.

\section*{Three editorial promises of this volume}

\begin{enumerate}
\item \textbf{The closed sector is not isolated.}  Every statement about bulk algebras is framed together with its boundary and line avatars.
\item \textbf{The homotopy upgrade is built in.}  Whenever a strict Poisson vertex algebra, dg algebra, or dg Lie algebra appears, the introduction points to the underlying $A_\infty$, brace, or $L_\infty$ package that actually governs it.
\item \textbf{Examples carry the theory.}  Free theories, Landau--Ginzburg models, $W$-algebras, Yang--Mills, and holographic systems are treated as the places where the general formalism becomes legible.
\end{enumerate}

\section*{Chapter-scale introductions}

\subsection*{Introduction}
The opening chapter should be read as a declaration that the bar complex has crossed a threshold.  What was a logarithm in Volume~I is now a Swiss-cheese algebraic object with an open and a closed side, and the reader is invited to follow that object into holomorphic-topological quantum field theory.  The introduction should therefore present not only the algebraic claims but also the physical image: local bulk operators, boundary operators, line operators, and their configuration-space OPE's all arranged around one common operadic skeleton.

\subsection*{The bar complex as Swiss-cheese algebra}
This chapter must be introduced as the structural heart of the volume.  The bar complex is not merely compatible with Swiss-cheese operations; it is the natural closed/open object on which those operations live.  The reader should feel that the two-colored operad is not a decorative framework but the precise geometry of bulk and boundary collisions.

The homotopy message is essential.  Swiss-cheese structure never honestly stops at strict multiplication and comultiplication.  Nested collisions and mixed open/closed degenerations generate higher operations automatically.  This chapter should make that inevitability visible from the first paragraph.

\subsection*{Holomorphic-topological locality}
This chapter should be introduced as the chapter that teaches the reader what sort of locality is actually available in the 3d theories of interest.  One has holomorphic dependence in one direction, topological invariance in another, and a resulting OPE calculus that is neither purely chiral nor purely topological.  The volume's operator algebras live precisely in this mixed locality.

The right introduction therefore emphasizes that every later algebraic structure is constrained by this hybrid geometry.  In particular, the higher operations are not arbitrary enrichments; they are what locality becomes when one cannot compress the theory back to a purely strict operator algebra.

\subsection*{Concrete axioms for $A_\infty$ chiral algebras}
This chapter should be introduced as the moment where the volume stops hinting at higher coherence and writes it down.  The aim is not to multiply axioms but to make explicit the minimal higher package needed to treat 3d HT theories honestly.  The reader should see these axioms as a carefully extracted normal form of what configuration spaces and BRST consistency already know.

A good introduction should stress that the reward is conceptual clarity: once the $A_\infty$ chiral language is accepted, the deformation-theoretic, physical, and geometric pieces of the volume all begin speaking the same dialect.

\subsection*{From $W(\mathsf{SC}^{\mathrm{ch,top}})$ to axioms and back}
This chapter should be introduced as the equivalence chapter in the strongest sense.  It proves that the abstract operadic packaging and the concrete axiomatics are not rival descriptions but mutually reconstructing views.  The reader should feel that the chapter settles the question of whether the formalism is too abstract by showing that nothing is lost in passing between pictures.

The homotopy upgrade appears here as a translation principle: what looks like an unwieldy operadic action in one language becomes an explicit hierarchy of products and homotopies in the other, and vice versa.  This is the bridge that lets the rest of the volume alternate freely between geometry and formulas.

\subsection*{BV-BRST construction}
This chapter should be introduced as the point where the Swiss-cheese story is generated directly from field theory.  The BV-BRST formalism does not merely supply examples; it explains why the entire higher package exists, how anomalies enter, and why deformation problems naturally become Maurer--Cartan equations.

The introduction should prepare the reader to see every later $A_\infty$ operation as a regularized remnant of a BRST anomaly calculation.  That physical meaning is what keeps the formalism honest.

\subsection*{Fulton--MacPherson calculus and proof of $A_\infty$ relations}
This is the geometric proof chapter.  It should be introduced as the place where the higher associativity relations are not asserted but read off from compactified configuration spaces and their boundary decompositions.  The reader should feel that every $A_\infty$ identity is a Stokes theorem waiting to be recognized.

The great virtue of this chapter is that it restores visibility.  Higher operations often look arbitrary in algebraic notation; here they are edges, vertices, collisions, and boundary faces.  This is one of the volume's clearest moments of ``show, don't tell.''

\subsection*{Cohomological descent: from $A_\infty$ to Poisson vertex algebras}
This chapter should be introduced as the first controlled shadow map of the volume.  Starting with a genuinely homotopy-coherent chiral structure, one descends to cohomology and obtains a shifted Poisson vertex algebra.  The point is not merely to recover a familiar object, but to show exactly what information is retained and what is compressed.

This chapter also delivers one of the cleanest explanations of why earlier quadratic and dg treatments are partial truths.  They are the cohomological silhouettes of a much richer chain-level world.

\subsection*{From the $A_\infty$-chiral structure to a Poisson vertex algebra on cohomology}
Where the previous chapter gives the descent mechanism, this one should be introduced as the sharpening of the resulting algebraic image.  It explains how the product, $\lambda$-bracket, and associated identities are distilled from the higher open/closed package.

The introduction should alert the reader that this is an act of controlled forgetting.  The PVA is powerful and computable, but the volume's philosophy is that it must always be read together with the chain-level source from which it descended.

\subsection*{Chiral Hochschild cohomology and bulk--boundary correspondence}
This chapter should be introduced as the first explicit articulation of the open/closed principle in the present setting.  Bulk operators are reconstructed via Hochschild-type objects, and the boundary begins to look like the open sector whose derived center remembers the closed theory.

This is one of the places where the volume most directly echoes the B-model, yet the introduction should stress that the holomorphic-topological setting is richer: higher operations, locality constraints, and spectral structures all intervene.

\subsection*{Chiral Bar--Cobar Duality and Koszul Duality Beyond Quadratic}
This chapter should be introduced as the explicit answer to one of the user's main demands.  The quadratic world is too small for Virasoro, for principal $\mathcal W$-algebras, for filtered boundary theories, and for many physically natural examples.  The purpose here is to show that bar--cobar duality and Koszul reasoning survive beyond the quadratic locus once one works at the correct homotopical level.

The introduction should make the reader feel the release: we are no longer trapped by presentations with only quadratic relations.  OPE's of arbitrary pole order, curved corrections, and modular data are absorbed into a single enlarged bar-cobar language.

\subsection*{Chiral Hochschild cochains as a brace algebra on $\FM(\C)\times\Conf(\R)$}
This chapter should be introduced as the moment where Hochschild theory becomes visibly operadic again.  Brace operations are not formal gadgets attached after the fact; they are the natural interaction algebra of local observables on mixed configuration spaces.  The reader should see the brace algebra as the grammar of iterated insertion.

This is also where the volume aligns most clearly with the recent higher-operations literature.  The higher brackets controlling deformations, generalized OPE's, and anomalies are not parallel to the brace story; they are its physical realization.

\subsection*{Line Operators, Koszul Duality, and Spectral $R$-Matrices}
This chapter should be introduced as the entrance of the line world proper.  Objects, morphisms, tensor products, and OPE's of lines force one to categorify the algebraic story.  The reader should feel that the appearance of spectral $R$-matrices is not a surprise but a natural consequence of moving from local operators to line operators in a mixed holomorphic-topological setting.

The homotopy upgrade is again unavoidable.  The category of lines is an $A_\infty$ world, its endomorphism algebra is a Koszul dual object with higher structure, and the $R$-matrix is part of that coherent package, not merely a strict endomorphism.

\subsection*{Bulk--boundary functoriality and spectral $R(z)$}
This chapter should be introduced as the dynamical refinement of the previous one.  Now the spectral parameter is not just present; it moves the tensor product, controls braiding, and records how bulk and boundary information is transported along the line category.  The reader should understand that $R(z)$ is a propagation law.

The introduction should also tell the reader that the Yang--Baxter equation is no longer a detached algebraic identity.  It is a coherence law for bulk--boundary-line compatibility.

\subsection*{The bulk--boundary--line Koszul triangle}
This chapter is the conceptual apex of the volume.  It should be introduced as the statement that the closed, open, and one-dimensional sectors are not three stories but three corners of a single duality triangle.  The reader should feel the geometry of the triangle before seeing its formulas: bulk centers, boundary categories, and line endomorphism algebras reflecting one another.

This is also where the full homotopy posture pays off most cleanly.  The triangle becomes false or badly distorted if any corner is artificially strictified too early.  The introduction should say that explicitly.

\subsection*{Celestial and boundary transfer: wedge algebras, obstruction towers, and the Airy--Witt realization}
This chapter should be introduced as a transport theorem chapter.  The same bar-cobar and boundary formalism that organizes the holomorphic-topological twist is now asked to move between boundary, celestial, and wedge-algebra settings.  The reader should experience this as a change of coordinates rather than a change of subject.

The right introduction should also prepare the reader for obstruction towers: transfer is never just functorial transport of a strict algebra, but transport of a homotopy package through a filtration whose failure modes are themselves meaningful invariants.

\subsection*{Physical origins}
This chapter should be introduced as the sourcebook for the entire volume.  Twists of 3d $\mathcal N=2$ theories, defect systems, holographic setups, and Poisson vertex/gauge-theory constructions are all contributing physical scenarios from which the algebraic objects of the volume emerge.  The point is not to justify the mathematics by physics but to show why the mathematics had to grow in this direction.

\subsection*{Holomorphic-topological field theories}
This chapter should be introduced as the chapter that teaches the reader how the mixed-dimensional field theories themselves are organized: what the superfields are, how the propagators look, what the topological and holomorphic directions do, and why factorization structures arise from them.  This is the habitat chapter.

A Chriss--Ginzburgified introduction should make the theory feel spatial and kinetic before it becomes formal.  The fields move; the propagators stretch; defects sit; brackets are produced.  Then the algebraic packaging becomes natural.

\subsection*{Holomorphic-topological boundary conditions and 4d origins}
This chapter should be introduced as the chapter where the volume's boundary sector is tied back to higher-dimensional parent theories.  Boundary conditions are not invented ad hoc in three dimensions; they are inherited, reduced, and twisted from richer worlds.  The reader should feel the genealogy.

\subsection*{The Heisenberg Rosetta Stone}
This chapter should be introduced as the model example of Volume~II.  Just as the Heisenberg frame opened Volume~I, the Rosetta Stone translates every new object here -- Swiss-cheese operations, brace structure, bulk-boundary maps, line modules, and spectral behavior -- into the simplest concrete language.  It is the chapter where the reader learns to read the new script fluently.

\subsection*{Explicit examples}
This chapter should be introduced as the general examples chapter in the best sense: not a miscellany, but a gallery arranged to display the range of the formalism.  The introduction should tell the reader what each example is testing: vanishing of higher products, emergence of a nontrivial $m_3$, spectral braiding, or transport across dualities.

\subsection*{Explicit Examples: Complete Computations of $A_\infty$ Operations}
This chapter is where the higher operations are fully performed rather than merely named.  The introduction should tell the reader that the reward for accepting the abstract formalism is that one can now compute nontrivial $A_\infty$ products exactly and watch the theory behave.

\subsection*{Worked examples}
This chapter should be introduced as the pedagogical mirror of the previous one.  Where the complete-computation chapter emphasizes exhaustiveness, the worked-examples chapter emphasizes intelligibility: how to move from a given Lagrangian or algebraic presentation to the actual operations and invariants.

\subsection*{$W$-Algebras as $A_\infty$ Chiral Algebras: Virasoro and $W_3$}
This chapter should be introduced as the decisive nonlinear test of Volume~II.  Virasoro and $W_3$ are precisely the examples that punish any attempt to remain in a strictly quadratic or strictly dg world.  The introduction should therefore announce them as the proving ground where the volume's full homotopy rhetoric must cash out in explicit structure.

\subsection*{Modular PVA quantization}
This chapter should be introduced as the frontier of quantization sharpened by modular data.  The question is not only how to quantize a classical PVA, but how to do so in a way that remembers boundary, genus, half-space, and planted-forest corrections.  The reader should feel that quantization here is a refined, geometric lifting problem.

\subsection*{The affine half-space BV theorem and its transport to Virasoro and $W_3$}
This chapter should be introduced as the theorematic heart of the quantization part.  The half-space is the arena in which affine current algebra, boundary chiral algebra, and BV transport are all simultaneously visible, and the subsequent transport to Virasoro and $W_3$ shows that the theorem is not local trivia but a mechanism.

\subsection*{The governing mechanism (FM$_3$ planted-forest synthesis)}
This chapter should be introduced as the combinatorial revelation of the volume.  The planted-forest picture explains why certain obstruction towers and modular corrections have the shape they do.  The reader should feel that a new geometry of organization has emerged, one that was latent in configuration spaces all along.

\subsection*{Toward the platonic ideal: twisted and untwisted Yang--Mills through bar--cobar, higher operations, instantons, and the infrared}
This chapter is where the user's requested tone should be felt most strongly.  It is the chapter of synthesis and ambition: the twisted and untwisted theories are placed into one landscape, the algebraic and physical constructions are made to answer each other, and the obstruction to a full infrared story is located with precision rather than handwaving.

The introduction should present Yang--Mills not as a remote application but as one of the true destinations of the whole bar-cobar programme.

\subsection*{Modular twisted/celestial holography from first principles}
This chapter should be introduced as the chapter where holography is no longer slogan but mechanism.  The reader should see how boundary chiral algebras, bulk observables, and Koszul duality package celestial and twisted holography into one algebraic architecture.  The chapter's core promise is that boundary algebras are not shadows after the fact; they are one of the principal engines of reconstruction.

\subsection*{The local law, the global law, and the bridge (logarithmic HT monodromy)}
This chapter should be introduced as the monodromy chapter: local collision data and global transport are shown to belong to the same law.  The reader should be told at once that the real protagonist is the bridge between them, not either side in isolation.

\subsection*{Anomaly-completed topological holography from first principles}
This chapter should be introduced as the place where the theory honestly confronts anomalies rather than treating them as nuisance terms.  The completion by anomalies is what makes the holographic correspondence structurally stable.  The reader should feel that this is not patching a defect in the theory but revealing one of its deepest organizing principles.

\subsection*{Conclusion and outlook}
The conclusion should be introduced not as a summary but as a redirection of vision.  The volume has shown that the Swiss-cheese/bar-cobar/PVA/line-operator world is coherent, computable, and structurally rich.  The outlook should explicitly point toward the remaining frontiers: nonperturbative line categories, full modularization, sharper raviolo comparisons, and geometric/holographic presentations of the resulting algebras.

\subsection*{Concordance: status ledger and cross-volume bridges}
This final unit should be introduced as the reader's navigational atlas.  It clarifies what is proved, what is conditional, what is conjectural, and exactly how Volume~II leans on Volume~I.  In a project of this size, concordance is itself a mathematical virtue: it keeps the architecture honest.

\subsection*{Relation to raviolo vertex algebras}
This appendix-scale unit should be introduced as a comparison in ontology.  Raviolo vertex algebras are not extraneous cousins but nearby incarnations of the same open/closed and local-to-global ideas.  The reader should understand what is gained and lost when passing between the Swiss-cheese and raviolo languages.

\subsection*{Raviolo restriction: Čech/Thom--Sullivan model and coinvariants}
The final unit should be introduced as a precise technical bridge.  Restriction, Čech descent, Thom--Sullivan models, and coinvariants provide the machinery for comparing local and global models in a language adapted to raviolo structures.  It is the sort of appendix that quietly guarantees that the bolder claims of the main text have a stable foundation.

\section*{Coda for Volume II}
Volume~II should now be read as the realization chapter of the whole monograph series.  What Volume~I built as algebraic engine and modular tower is here turned into a visible ecology of bulk, boundary, and line structures.  The core philosophical shift is this: the full homotopy package is no longer a correction to the story; it is the story.  Every strict object that survives in the text survives because a much larger coherent object has chosen, in that example, to cast a simple shadow.
